\section{Conclusion}

The ViPragSent resource provides a Vietnamese dataset for six pragmatic phenomena, and the proposed XLM-R-large multi-task model uses separate polarity, emotion, and training-only rationale objectives. On the 2,000-example test cohort, it obtains $93.67 \pm 0.19$ macro-pragmatic F1 and the highest observed mean on all six pragmatic heads among the evaluated systems with complete three-run results. Q1b--Q4 characterize external affective transfer, component and computational trade-offs, positive-label scarcity, and probability calibration, respectively. Rationale supervision is training-only and is not used by the inference interface.
